\documentclass{article}
\usepackage{tikz}
\usetikzlibrary{arrows.meta, positioning}

\usepackage[utf8]{inputenc}
\usepackage[spanish,mexico]{babel}
\usepackage{listings}
\usepackage{amsmath}
\setlength{\textwidth}{18cm}
\setlength{\oddsidemargin}{-1cm}
\setlength{\headsep}{1cm}
\setlength{\voffset}{0cm}
\setlength{\topmargin}{0cm}
\setlength{\headheight}{0cm}
\usepackage{tikz}
\usepackage{semantic}
\usepackage{url}
\usetikzlibrary{positioning}
\usetikzlibrary{calc,arrows}
\usepackage{multicol}
\usepackage{lipsum} 
\usepackage{multirow}
\usepackage{amssymb}

\usepackage{graphicx}
\usepackage{forest}
\usepackage{tikz-qtree}
\usepackage{xcolor}

\begin{document}
\pagecolor{black}
\color{white}

%%%%%% ENCABEZADO %%%%%%%%%%%%%%%%%%%%%%%%%%%%%%%%%%%%%%%
    \colorbox{black}{
        \begin{minipage}[t]{0.16 \textwidth}
           \begin{flushright}
            \includegraphics[width=1in]{UNAM.png}
           \end{flushright}
        \end{minipage}
        \begin{minipage}[H]{0.62 \textwidth}
            \begin{center}
                {\large \textsc{Universdad Nacional Autónoma de México}}
                \vspace{0.25cm}
                \\
                { \large \textbf{Lenguajes de Programacion\\ Examen Parcial IV}}                
                \textbf{}
                \begin{multicols}{2}
                \begin{flushleft}
                \begin{itemize}
                    % NOMBRES DE INTEGRANTES
                    \item  \small Edgar Montiel Ledesma\\ 317317794
    
                    \item \footnotesize Carlos Daniel Cortes Jimenez\\ 420004846
                \end{itemize}
                \end{flushleft}
                \vspace{0.25cm}
                \end{multicols} 
            \end{center}
            \vspace{0.05cm}
        \end{minipage}
        \begin{minipage}[t]{0.16 \textwidth}
            \begin{flushleft}
                \includegraphics[width=1in]{EFC.png}
            \end{flushleft}
        \end{minipage}
    }
    
    \begin{tikzpicture}
        \draw[thick] (-6.5,0)--(11.2,0);
    \end{tikzpicture}
    %%%%%%%%%%%%%%%%%%%%%%%%%%%%%%%%%%%%%%%%%%%%%%%%%%%%%%%%%
    \begin{itemize}
        \item[1.] Considera el siguiente programa en el lenguaje While :
        \begin{itemize}
            \item[ ] new z = 0;
            \item[ ] \, while (y < x + 1) do
            \item[ ] \, (z := z + 1;
            \item[ ] \, x := x - y)
            \item[ ] end
        \end{itemize}

        \begin{itemize}
            \item[a)] Ejecuta el programa en el estado en el que $\sigma$(x) = 17 y $\sigma$(y) = 5 ¿Cual es el estado resultante de la evaluacion?
            \item[b)] Da un estado $\sigma$ tal que si se evalua el programa anterior con dicho estado la evaluacion se ciclarıa infinitamente.
        \end{itemize}

        \item[2.] Extiende el lenguaje While con el operador:
        \begin{center}
        for x := a1 to a2 do c
        \end{center}
        esto es:
        \begin{itemize}
            \item[a)] Modifica la estructura de la maquina W (agregando marcos, estados o transiciones) para evaluar la expresion for.
            \item[b)] Da las reglas de semantica estatica para verificacion de tipos para el nuevo operador for.
            \item[c)] ¿Es posible definir el operador for como azucar sintactica dentro del lenguaje While? justifica tu respuesta.
        \end{itemize}

        \item[3.] Decimos que dos programas en el lenguaje While son equivalentes (c1 $\equiv$w c2) si y solo sı la ejecucion de ambos programas resulta en el mismo estado, es decir, si para todo estado de las variables $\sigma$, $\blacklozenge$ $\succ$  ⟨c1, $\sigma$⟩ $\rightarrow{*}$ W $\blacklozenge$ $\prec$ $\sigma'$ y $\blacklozenge$ $\succ$ ⟨c2, $\sigma$⟩ $\longrightarrow*W$ $\blacklozenge$ $\prec$ $\sigma'$ entonces c1 $\equiv$w c2.\\
        Con la definicion de equivalencia anterior, demuestra o da un contraejemplo de lo siguiente:
        \begin{itemize}
            \item[a)] $\equiv$w realmente es una relacion de equivalencia. Esto es, demuestra que la relacion $\equiv$w es transitiva, reflexiva y simetrica.
            \item[b)] c;skip $\equiv$w c
            \item[c)] c1; c2 $\equiv$w c2; c1
            \item[d)] c1; (c2; c3) $\equiv$w (c1; c2); c3 
        \end{itemize}
    \end{itemize}

\end{document}
